\section{Алгоритм Похлига--Хеллмана}

Алгоритм Похлига--Хеллмана предназначен для решения задачи дискретного логарифмирования в циклической группе в тех случаях, когда порядок рассматриваемой подгруппы имеет нетривиальное разложение на простые множители. Для задачи на эллиптических кривых это означает, что метод применим к подгруппе $\langle P\rangle$, если порядок точки $P$ известен и представим в виде
\[
n=\prod_{i=1}^{r} p_i^{e_i},
\]
где $p_i$ --- различные простые числа, а $e_i\ge 1$. Основная идея алгоритма состоит в том, чтобы не искать сразу весь дискретный логарифм $k$ из равенства $Q=kP$, а поочерёдно восстанавливать его остатки по модулям $p_i^{e_i}$, после чего объединять их в одно решение по модулю $n$.

Пусть требуется найти число $k=\log_P Q$. Тогда для каждого сомножителя $p_i^{e_i}$ вычисляется значение
\[
k_i \equiv k \pmod {p_i^{e_i}}.
\]
После этого получается система сравнений
\[
k \equiv k_1 \pmod {p_1^{e_1}}, \qquad
k \equiv k_2 \pmod {p_2^{e_2}}, \qquad \dots, \qquad
k \equiv k_r \pmod {p_r^{e_r}},
\]
которая имеет единственное решение по модулю $n$ в силу китайской теоремы об остатках. Тем самым исходная задача сводится к решению нескольких более простых задач, соответствующих простым степеням в разложении порядка группы.

Рассмотрим, как вычисляется одно из чисел $k_i$. Для краткости положим $p=p_i$ и $e=e_i$. Тогда число $k_i$ можно записать в $p$-ичной форме
\[
k_i=z_0+z_1p+z_2p^2+\dots+z_{e-1}p^{e-1},
\]
где каждая цифра $z_j$ принадлежит множеству $\{0,1,\dots,p-1\}$. Алгоритм восстанавливает эти цифры последовательно, начиная с младшего разряда. Для этого строится точка
\[
P_0=\frac{n}{p}P,
\]
имеющая порядок $p$. Далее по исходной точке $Q$ и уже найденным цифрам формируются новые точки, дискретные логарифмы которых вычисляются уже в подгруппе порядка $p$. Таким образом, вычисление $k_i$ по модулю $p^e$ сводится к последовательности из $e$ задач дискретного логарифмирования в существенно меньшей подгруппе порядка $p$.

Первый шаг имеет особенно простой вид. Полагая
\[
Q_0=\frac{n}{p}Q,
\]
получаем
\[
Q_0=\frac{n}{p}Q=\frac{n}{p}(kP)=k\frac{n}{p}P=kP_0.
\]
Так как точка $P_0$ имеет порядок $p$, число $k$ в этом выражении фактически рассматривается по модулю $p$, а значит,
\[
Q_0=z_0P_0.
\]
Следовательно, первая цифра $z_0$ находится как дискретный логарифм точки $Q_0$ по основанию $P_0$ в подгруппе порядка $p$.

После нахождения $z_0$ можно исключить вклад младшего разряда и перейти к следующему. Для этого рассматривается точка
\[
Q_1=\frac{n}{p^2}(Q-z_0P).
\]
Подстановка разложения $k_i=z_0+z_1p+\dots$ показывает, что
\[
Q_1=z_1P_0.
\]
Аналогично, если уже найдены цифры $z_0,z_1,\dots,z_{t-1}$, то определяется точка
\[
Q_t=\frac{n}{p^{t+1}}\left(Q-z_0P-z_1pP-\dots-z_{t-1}p^{t-1}P\right),
\]
и очередная цифра находится из равенства
\[
Q_t=z_tP_0.
\]
Тем самым вычисление дискретного логарифма по модулю $p^e$ сводится к последовательному определению коэффициентов $z_0,z_1,\dots,z_{e-1}$.

С алгоритмической точки зрения важнейшее достоинство метода заключается в том, что вся сложность переносится на решение задач в подгруппах малых простых порядков. Если простые множители $p_i$ невелики, то каждая такая подзадача решается сравнительно быстро, а итоговый дискретный логарифм восстанавливается объединением частичных результатов. Поэтому трудоёмкость алгоритма определяется не столько величиной всего порядка $n$, сколько размерами наибольших простых делителей в его разложении.

Именно отсюда вытекает криптографический смысл алгоритма Похлига--Хеллмана. Если порядок точки-генератора раскладывается на малые простые множители, то задача дискретного логарифмирования перестаёт быть задачей большого размера и распадается на набор существенно более простых задач. Следовательно, подгруппа, используемая в криптографической схеме, должна иметь большой простой порядок либо, по крайней мере, содержать большой простой делитель. Иначе стойкость схемы резко снижается.